\section*{Overview}

The number theoretic transform is the modular counterpart of FFT. In contest terms, it is the standard exact method
for multiplying polynomials in \(O(n \log n)\) time when the coefficients live modulo a suitable prime.

Whenever a counting problem secretly asks for convolution, NTT is often the cleanest fast solution if the modulus fits.

\section*{When to Use It}

Use it when:

\begin{itemize}
  \item you need polynomial convolution faster than \(O(n^2)\),
  \item coefficients are naturally taken modulo a prime such as \(998244353\),
  \item exact modular arithmetic is preferable to floating-point FFT,
  \item the problem can be reframed as counting pair sums, subset combinations, or polynomial products.
\end{itemize}

For tiny input sizes, NTT is overkill. For incompatible moduli, a different transform strategy or CRT layer may be
needed.

\section*{Core Idea}

Convolution becomes pointwise multiplication after transforming to the right basis. FFT uses complex roots of unity.
NTT does the same thing inside modular arithmetic using roots of unity modulo a prime.

For the common modulus \(998244353\),
\[
998244353 - 1 = 119 \cdot 2^{23},
\]
so there are enough \(2^k\)-th roots of unity for transform lengths up to \(2^{23}\). That is exactly why this modulus
appears so often in contests.

\section*{Key Insight}

The transform is fast because the evaluation points are highly structured. Instead of evaluating a polynomial at every
point independently, the butterfly layers reuse work between even and odd parts.

The other key fact to remember is that the inverse transform is not ``free''. After reversing the butterfly process,
every coefficient must be multiplied by \(n^{-1} \bmod \texttt{MOD}\).

\section*{Operations / Main Technique}

The standard convolution pipeline is:

\begin{itemize}
  \item choose a power-of-two length \(n\) large enough for the result,
  \item pad both arrays to length \(n\),
  \item run the forward transform on both,
  \item multiply pointwise,
  \item run the inverse transform,
  \item normalize by \(n^{-1}\).
\end{itemize}

\paragraph{Worked viewpoint.}
If \(A(x)\) and \(B(x)\) are transformed into their values at the same roots of unity, then
\[
(A \cdot B)(\omega_i) = A(\omega_i) B(\omega_i).
\]
That is the whole reason pointwise multiplication in transform space corresponds to convolution in coefficient space.

\section*{Correctness Intuition}

The chosen roots of unity make the transform matrix invertible, just like the complex DFT. So the forward transform is
an evaluation map at special points, and the inverse transform reconstructs the unique polynomial coefficients from
those values.

The butterfly layers are only a structured implementation of that algebra. Pointwise multiplication after the forward
transform gives the values of the product polynomial at those same points, and the inverse transform recovers the
coefficients.

\section*{Complexity Analysis}

\begin{itemize}
  \item forward transform: \(O(n \log n)\),
  \item inverse transform: \(O(n \log n)\),
  \item convolution: \(O(n \log n)\) total after padding,
  \item memory: \(O(n)\).
\end{itemize}

Here \(n\) is the padded transform length, not the original polynomial degree.

\section*{Implementation}

The reference code uses:

\begin{itemize}
  \item modulus \(998244353\),
  \item primitive root \(3\),
  \item iterative bit reversal,
  \item a \texttt{convolution(a, b)} wrapper.
\end{itemize}

That version is the contest-ready baseline I want to understand before hiding everything behind a library call.

\section*{Common Pitfalls}

\begin{itemize}
  \item Using a modulus that does not support the required transform length.
  \item Forgetting inverse normalization by \(n^{-1}\).
  \item Padding to a length smaller than \(|a| + |b| - 1\).
  \item Mixing negative values into modular arithmetic without re-normalizing.
  \item Using NTT when a quadratic convolution would have been simpler and fast enough.
\end{itemize}

\section*{Variants / Extensions}

\begin{itemize}
  \item CRT-based convolution for arbitrary moduli,
  \item AtCoder Library convolution wrappers,
  \item formal power series routines built on repeated NTT,
  \item subset of problems where FFT with doubles is acceptable instead.
\end{itemize}

For most CP uses, direct polynomial multiplication is only the entry point. Once NTT is available, many generating
function and counting tricks become feasible.

\section*{Practice Problems}

\begin{itemize}
  \item Multiply two polynomials modulo \(998244353\).
  \item Count pair sums by multiplying frequency polynomials.
  \item DP transitions that become convolution after regrouping states.
\end{itemize}

\section*{References}

\begin{itemize}
  \item \href{https://cp-algorithms.com/algebra/fft.html}{cp-algorithms: FFT and NTT overview}.
  \item \href{https://atcoder.github.io/ac-library/production/document_en/convolution.html}{AtCoder Library: convolution}.
  \item \href{https://codeforces.com/catalog}{Codeforces Catalog}.
\end{itemize}
